\chapter{Beyond Basic Graph Patterns}
\label{ch:sparql}

A basic graph pattern is a conjunction: every triple pattern must match, or the
row is gone. Most questions people ask of an ontology are not conjunctions.
``List every material, with its melting point \emph{if we know it}'' keeps a
material whose melting point is missing. ``Metals or ceramics'' is a
disjunction. ``How many subclasses does each class have'' collapses rows into
groups. This chapter adds the operators behind those questions, \code{OPTIONAL},
\code{UNION}, and aggregation, and the update form whose triples come from a
query. Each one is a relational operator the earlier chapters did not need,
and each one constrains the optimizer in a way an inner join does not.

\section{Where a conjunction loses rows}

Take three materials, two of them with a melting point:
\begin{lstlisting}
:Steel   :meltsAt 1370 .
:Copper  :meltsAt 1085 .
:Steel a :Material .  :Copper a :Material .  :Glass a :Material .
\end{lstlisting}
The pattern \code{\{ ?m a :Material . ?m :meltsAt ?t \}} is a join of two
scans on \code{?m} (\cref{ch:joins}), and it returns two rows. \code{Glass} is
not wrong; it is unknown, and the inner join cannot say ``unknown''. The fix is
an operator whose output keeps every row of one side.

\section{OPTIONAL is a left outer join}

SPARQL's algebra works on \term{solution mappings}: partial functions from
variables to terms. Two mappings are \term{compatible} when they agree on every
variable both bind. For multisets of mappings $\Omega_1$ and $\Omega_2$ and a
filter expression $e$ \cite{sparql,perez2009}:
\[
\mathit{LeftJoin}(\Omega_1, \Omega_2, e) \;=\;
\mathit{Filter}\big(e, \mathit{Join}(\Omega_1, \Omega_2)\big) \;\cup\;
\{\, \mu_1 \in \Omega_1 \mid \neg\exists \mu_2 \in \Omega_2 :\ \mu_1 \sim \mu_2 \wedge e(\mu_1 \cup \mu_2) \,\}.
\]
The first part is the inner join; the second keeps each left mapping that found
no compatible partner, with the right side's variables \term{unbound}. For the
example, \code{?m a :Material OPTIONAL \{ ?m :meltsAt ?t \}} gives:

\begin{center}
\begin{tabular}{@{}ll@{}}
\toprule
\code{?m} & \code{?t} \\
\midrule
\code{:Steel} & 1370 \\
\code{:Copper} & 1085 \\
\code{:Glass} & \emph{(unbound)} \\
\bottomrule
\end{tabular}
\end{center}

An unbound value is not a value. It is a hole in the row, and the executor has
to carry it: rows gain a ``not bound'' marker per column, which the result
printer shows as an empty cell.

\subsection{A filter inside OPTIONAL is part of the join}

The $e$ in the definition matters. Compare
\begin{lstlisting}
SELECT ?m ?t WHERE { ?m a :Material OPTIONAL { ?m :meltsAt ?t FILTER(?t > 1200) } }
SELECT ?m ?t WHERE { ?m a :Material OPTIONAL { ?m :meltsAt ?t } FILTER(?t > 1200) }
\end{lstlisting}
The first keeps all three materials: \code{Steel} with 1370, and \code{Copper}
and \code{Glass} with \code{?t} unbound, because the filter only decides which
right rows may join. The second filters the finished rows, and a comparison
with an unbound value is an \emph{error}, which \code{FILTER} treats as false: it
returns \code{Steel} alone. Moving a filter across an \code{OPTIONAL} boundary
changes the answer, so the planner must not do it.

\begin{keyidea}
An outer join is not commutative, and filters do not move freely across it. The
left side is preserved; the right side only decorates it.
\end{keyidea}

\section{Executing a left join}

Both join algorithms of \cref{ch:joins} extend to outer joins with one extra bit
of state.

\begin{description}[leftmargin=1.2em]
\item[Nested loop.] For each left row, rebind and scan the right side. If no
right row matched (after the join filter), emit the left row once with the
right columns unbound. The bit is ``did this left row match'', and it lives for
one iteration of the outer loop.
\item[Hash join.] Build the hash table on the \emph{right} side and probe with
the left. A probe that finds no partner emits the left row padded. If the build
side were the left one, every left row would need a ``matched'' flag in the
table and a final pass over the table to emit the unmatched ones; that works,
but it makes the operator blocking at its end as well as its start.
\end{description}

An index nested-loop left join is the common case in practice: the left side
is a selective pattern, and the optional part is a lookup by subject on SPO.

\begin{decision}{Left joins build on the right (planned)}
\dpart{Choice} The hash left join builds on the optional (right) side and
probes with the preserved (left) side.
\dpart{Alternatives} Build on the left with a matched flag per entry, and emit
unmatched entries after the probe phase.
\dpart{Why} Padding happens at probe time, the operator streams its output,
and the optimizer's choice is only whether the right side fits in memory.
\dpart{Cost} When the right side is much larger than the left, building on it
costs memory a left build would not.
\end{decision}

\section{UNION}

\code{\{ P1 \} UNION \{ P2 \}} evaluates both sides and concatenates their rows.
SPARQL union is a bag union: a row produced by both sides appears twice, and
\code{DISTINCT} removes duplicates across sides. The sides may bind different
variables, so a row from one side has the other side's variables unbound, the
same marker \code{OPTIONAL} introduced. The executor is the simplest operator in
the engine, a concatenation of two child iterators; the work is in the binder,
which computes the output columns as the union of both sides' variables.

\section{Aggregation}

\code{GROUP BY} partitions rows by the values of its keys and computes one row
per group:
\begin{lstlisting}
SELECT ?c (COUNT(?sub) AS ?n) WHERE { ?sub rdfs:subClassOf ?c }
GROUP BY ?c HAVING (COUNT(?sub) > 2)
\end{lstlisting}
There are two classic implementations \cite{graefe}:
\begin{description}[leftmargin=1.2em]
\item[Sort aggregation.] Sort by the group keys (the external merge sort of
\cref{ch:joins} already exists), then scan once, finishing a group whenever the
key changes. Output comes out ordered by the keys.
\item[Hash aggregation.] Keep a hash table from group key to an accumulator and
update it per row. One pass, no sort, but the table must hold every group.
\end{description}
Accumulators are small: \code{COUNT} is a counter, \code{SUM}, \code{MIN}, and
\code{MAX} one value, \code{AVG} a sum and a count, divided at the end.
\code{COUNT(DISTINCT x)} is the exception: it needs the set of distinct values
per group, so its memory grows with the data, not with the number of groups.
\code{HAVING} is a filter on the output of the aggregation.

\begin{decision}{Hash aggregation, and what to do when groups outgrow memory (planned)}
\dpart{Choice} Hash aggregation. When the table exceeds its memory budget,
the plan records whether to spill or refuse; the leading design is to partition
the input by a hash of the key into files, then aggregate each partition in
memory, as hybrid hash join does for joins.
\dpart{Alternatives} Sort aggregation always; or refuse with a clear error past
the budget.
\dpart{Why} Ontology workloads group by class or property, and there are far
fewer of those than triples, so the table almost always fits and one pass
beats a sort.
\dpart{Cost} Output is unordered, so \code{ORDER BY} on the keys still sorts;
the spill path is code that runs rarely, which is why it needs its own tests.
\end{decision}

\section{Planning around the new operators}

The join-order search of \cref{ch:joins} treats a basic graph pattern as a set
of patterns to permute. With \code{OPTIONAL}, the pattern becomes a tree, and
only parts of it may be reordered:
\begin{itemize}
\item Patterns inside the left side may be reordered among themselves, and so
may patterns inside the right side. A pattern may not move from one side of an
\code{OPTIONAL} to the other: that would turn an optional match into a required
one, or the reverse.
\item The cardinality of a left join is at least that of its left side, since
every left row survives. An estimate below that is a bug, not a guess.
\item \code{\textbackslash explain analyze} shows the new operators with
estimated and actual rows, as it does for joins, so a bad estimate is visible.
\end{itemize}
Pérez, Arenas, and Gutierrez \cite{perez2009} identify the \term{well-designed}
patterns, those where every variable shared between an optional part and the
outside also appears on the left; for them, more rewrites are safe, and
evaluation is easier. Patterns written by people are almost always
well-designed.

\section{Updates whose triples come from a query}

\code{DELETE \{ \ldots \} INSERT \{ \ldots \} WHERE \{ \ldots \}} rewrites every
match: move every class under \code{:Metal} to \code{:MetallicMaterial},
for example. The order of the three parts is the whole semantics:

\begin{enumerate}
\item Evaluate the \code{WHERE} pattern \emph{in full}, producing every solution.
\item Instantiate and delete the \code{DELETE} templates for every solution.
\item Instantiate and insert the \code{INSERT} templates for every solution.
\end{enumerate}

Interleaving them is wrong for the same reason as the \emph{Halloween problem}
of relational updates: a scan that sees its own inserts can match them again,
and an update that adds a superclass link per subclass could loop forever or
apply twice. Materializing the solutions first gives the update a snapshot of
the data as it was before the statement, which is what the specification
means. The whole statement is one transaction, under the same two-phase locks
as \code{INSERT DATA} and \code{DELETE DATA} (\cref{ch:concurrency}).

\begin{decision}{Evaluate WHERE first, then delete, then insert (planned)}
\dpart{Choice} The solutions of \code{WHERE} are materialized before any triple
changes; deletes run before inserts.
\dpart{Alternatives} Pipeline solutions straight into the updates; or insert
before delete.
\dpart{Why} It is the specification's semantics and it avoids the Halloween
problem. Deleting first means a triple both deleted and inserted by one
statement ends up present, which is what a rewrite intends.
\dpart{Cost} The solutions are held in memory (or spilled) for the statement's
duration.
\end{decision}

\begin{status}
Planned (plan 12.1--12.5): today \code{OPTIONAL}, \code{UNION}, aggregates, and
\code{DELETE/INSERT WHERE} are parsed far enough to fail with a line and a
column. Every new form will run on both backends and join the differential test
that compares the indexed store with the memory store.
\end{status}

\begin{checkyourself}
\item Write the rows of \code{?m a :Material OPTIONAL \{ ?m :meltsAt ?t FILTER(?t < 1200) \}} for the three materials above.
\item Why can a hash left join not simply swap its build and probe sides when the left side is smaller, as an inner hash join can?
\item A union's left side binds \code{?a} and its right side binds \code{?b}. What columns does the result have, and what is in them?
\item Why does \code{COUNT(DISTINCT ?x)} need more memory than \code{COUNT(?x)}, and how does that interact with hash aggregation?
\item An optimizer estimates 40 rows for a left join whose left input it estimates at 100. What is wrong?
\item \code{INSERT \{ ?c rdfs:subClassOf :Thing \} WHERE \{ ?c rdfs:subClassOf ?d \}} runs with the inserts visible to the scan. What can go wrong, and how does materializing \code{WHERE} first prevent it?
\end{checkyourself}

\section*{Further reading}
The SPARQL 1.1 specification \cite{sparql}, section 18, defines the algebra;
Pérez, Arenas, and Gutierrez \cite{perez2009} give its semantics and
complexity, including well-designed patterns; Graefe's survey \cite{graefe}
covers outer joins and hash and sort aggregation.
